\documentclass[10pt,a4paper]{amsart}
\usepackage[margin=19mm]{geometry}
\usepackage{amsmath,amssymb,mathtools}
\usepackage[scr=rsfs,cal=euler]{mathalfa}
\usepackage{xfrac}
\usepackage[unicode,hidelinks,bookmarks=false]{hyperref}
\newcommand{\ie}{i.e.\ }
\newcommand{\cf}{cf.\ }
\newcommand{\resp}{respectively\ }
\newcommand{\Subsection}{\subsection}
\providecommand{\nobreakdash}{\nobreak\hbox{-}}
\setlength{\emergencystretch}{2em}
\multlinegap0pt
\usepackage{mathtools}
\usepackage{amssymb}
\usepackage[shortlabels]{enumitem}
\usepackage{bm}
\usepackage{float}
\usepackage{xcolor}
\usepackage{tikz}
\newtheorem{thm}{Theorem}
\newtheorem{prop}{Proposition}
\title{The Hopf decomposition of locally compact group actions}
\author{Nachi Avraham-Re'em, George Peterzil}
\date{Source: arXiv:2406.17137v3 (CC BY 4.0)}
\begin{document}
\maketitle

\noindent\emph{Source and licence.} The Hopf decomposition of locally compact group actions. Version arXiv:2406.17137v3 (CC BY 4.0), distributed by arXiv under the Creative Commons Attribution 4.0 International licence, http://creativecommons.org/licenses/by/4.0/, as recorded in the arXiv licence metadata for this exact version and shown on the abstract page https://arxiv.org/abs/2406.17137v3. Full licence notice: \texttt{licenses/arxiv-2406.17137-CC-BY-4.0.txt}.\par\smallskip

\noindent\textit{Editorial scope.} Counted unit: the article's thm thm (source0.tex lines 912--919). The article's complete original proof of it is delivered with it (source0.tex lines 921--956, one \texttt{proof} environment of 36 lines). No mathematics is written, repaired or re-derived: the statement and the proof are the article's own lines, in the article's own notation, and no retained line is substituted or re-flowed.  Retained uncounted as the source-local context the proof needs: lines 511--514; lines 1131--1134; lines 1313--1328; lines 1340--1348.  Nothing was omitted from any retained span, so the package carries no omission record.  No retained line contains a citation, so the wrapper carries no bibliography and no bibliography entry.  No retained line contains a figure environment, an image-inclusion command or drawing code, so no image is delivered and the package declares no asset dependency.  Editorial formatting only, in addition: an AMS \texttt{amsart} preamble that restates the article's own declarations and macros used by the retained lines, namely the environments \texttt{thm}, \texttt{prop} on separate counters, together with the standard packages those lines need.  The retained environments print in this document as Theorem 1, Proposition 1 in order of appearance, whereas the article prints the same items with the numbering of its own full document.  The complete native source of the article as distributed in the arXiv e-print for this exact version is bundled unchanged as \texttt{sources/161319/source0.tex}.
\begin{thm}[Following Krengel--Rosinski]
\label{thm:disssmooth}
If \(\left(X,\mu\right)\) is a totally dissipative nonsingular \(G\)-space, then \(E_{G}^{X}\) is essentially smooth.
\end{thm}

\begin{prop}
\label{prop:dissacim}
For every nonsingular \(G\)-space \(\left(X,\mu\right)\), the dissipative part \(\mathcal{D}\) admits an {\sc a.c.i.m} equivalent to \(\mu\mid_{\mathcal{D}}\).
\end{prop}

\begin{thm}[Measures on homogeneous spaces]
\label{thm:meashom}
For a group pair \(C\lessdot G\):
\begin{enumerate}
    \item There is a unique, up to a positive constant multiple, invariant \(\sigma\)-finite Radon measure \(\kappa\) on \(G/C\) iff \(\varDelta_{G}\mid_{C}\equiv\varDelta_{C}\).
    \item There is always a unique, up to measure class, quasi-invariant \(\sigma\)-finite Radon measure \(\kappa\) on \(G/C\), with continuous Radon--Nikodym cocycle.
    \item \emph{Weil's formula}: For every \(\kappa\) as in (2) there is a rho-function \(\rho\) with
    \begin{equation}
    \begin{aligned}
    \label{eq:weil1}
    \int_{G}\varphi\left(g\right)d\lambda_{G}\left(g\right)=\int_{G/C}\Big[\int_{C}\varphi\left(gc\right){\textstyle \frac{\varDelta_{G}\left(c\right)}{\varDelta_{C}\left(c\right)}}d\lambda_{C}\left(c\right)\Big]\rho\left(g\right)^{-1}d\kappa\left(gC\right),
    \end{aligned}
    \end{equation}
    \[\text{ for every Borel function } \varphi:G\to\mathbb{R}_{\geq 0}.\]
\end{enumerate}
\end{thm}

\begin{thm}[Mackey]
\label{thm:haaruniq}
For a group pair \(C\lessdot G\) with \(C\) compact, with the invariant measure \(\kappa\) constructed in Theorem~\ref{thm:meashom}(1), the following hold:
\begin{enumerate}
    \item Every quasi-invariant \(\sigma\)-finite Borel measure on \(G/C\) is equivalent to \(\kappa\).
    \item Every invariant \(\sigma\)-finite Borel measure on on \(G/C\) is a positive constant multiple of \(\kappa\).
    \item The pushforward of \(\lambda_{G}\) along the canonical projection \(G\to G/C\) is a positive constant multiple of \(\kappa\).
\end{enumerate}
\end{thm}

\begin{thm}
For every essentially free nonsingular \(G\)-space \(\left(X,\mu\right)\), the following are equivalent:
\begin{enumerate}
    \item \(\left(X,\mu\right)\) is totally dissipative.
    \item \(E_{G}^{X}\) is essentially smooth and \(\mu\) admits a \(G\)-invariant equivalent measure.
    \item \(\left(X,\mu\right)\) is isomorphic to a translation \(G\)-space.
\end{enumerate}
\end{thm}

\begin{proof}
(1)\(\implies\)(2). This follows directly from Theorem~\ref{thm:disssmooth} and Proposition~\ref{prop:dissacim}.

(2)\(\implies\)(3). Since \(E_{G}^{X}\) is essentially smooth it has a Borel transversal \(W_{o}\subseteq X\). Consider the restriction of the action map to \(W_{o}\), namely
\[\Phi:W_{o}\times G\to X,\quad \Phi\left(w,g\right)=g.w.\]
Using the transversal property of \(W_{o}\) and that the action of \(G\) is free, \(\Phi\) is injective Borel map, so it admits a Borel inverse that we denote
\[\Phi^{-1}\left(x\right)=\left(\omega\left(x\right),\gamma\left(x\right)\right)\text{ where }\omega:X\to W_{o}\text{ and }\gamma:X\to G.\]
We claim that
\[\omega\left(g.x\right)=\omega\left(x\right)\text{ and }\gamma\left(g.x\right)=g\gamma\left(x\right)\text{ for all }x\in X,\,\,g\in G.\]
Indeed, the first identity is by the transversal property, and this implies
\[g\gamma\left(x\right).\omega\left(x\right)=g.x=\gamma\left(g.x\right).\omega\left(g.x\right)=\gamma\left(g.x\right).\omega\left(x\right),\]
which implies the second identity since the action is free.

Fix some invariant measure \(\eta_{o}\) equivalent to \(\mu\). Consider the measure \(\mu_{1}\coloneqq \eta_{o}\circ\Phi\). Then for every Borel sets \(A\subseteq W_{o}\) and \(B\subseteq G\), and every \(g\in G\),
\begin{align*}
&\mu_{1}\left(A\times g^{-1}B\right)\\
&\quad=\eta_{o}\left(\Phi\left(A\times g^{-1}B\right)\right)\\
&\quad=\eta_{o}\left(x\in X:\omega\left(x\right)\in A,g\omega\left(x\right)\in B\right)\\
&\quad=\eta_{o}\left(x\in X:\omega\left(g.x\right)\in A,\omega\left(g.x\right)\in B\right)\\
&\quad=\eta_{o}\left(x\in X:\omega\left(x\right)\in A,\omega\left(x\right)\in B\right)\\
&=\mu_{1}\left(A\times B\right).
\end{align*}
It follows that for every Borel set \(A\subseteq W_{o}\), the map \(B\mapsto\mu_{o}\left(A\times B\right)\) is a left invariant Borel \(\sigma\)-finite measure on \(G\), so it is \(\lambda\) up to a positive constant (see Theorem~\ref{thm:haaruniq}(2)). Thus, there is a positive constant \(\mu_{o}\left(A\right)\) such that
\[\mu_{1}\left(A\times B\right)=\mu_{o}\left(A\right)\lambda\left(B\right).\]
As \(\mu_{1}\) is a measure, \(A\mapsto\mu_{o}\left(A\right)\) defines a measure on \(W_{o}\) so we deduce (3).

(3)\(\implies\)(1). For a translation \(G\)-space attached to \(\left(W_{o},\nu_{o}\right)\), since it is measure preserving we may assume that \(\nabla\left(w,h\right)=1\) for every \(\left(w,h\right)\in W_{o}\times G\), so for every \(f\in L_{+}^{1}\left(W_{o}\times G,\nu_{o}\otimes\lambda\right)\) we have
\[S_{f}^{G}\left(w,h\right)=\int_{G}f\left(w,gh\right)d\lambda\left(g\right).\]
By the Fubini Theorem and the integrability of \(f\) we have
\begin{align*}
\int_{W_{o}}S_{f}^{G}\left(w,h\right)d\nu_{o}\left(w\right)	
&=\iint\nolimits_{W_{o}\times G}f\left(w,gh\right)d\nu_{o}\otimes\lambda\left(w,h\right)\\
&=\iint\nolimits_{W_{o}\times G}f\left(w,h\right)d\nu_{o}\otimes\lambda\left(w,h\right)<+\infty.
\end{align*}
It follows that \(S_{f}^{G}\left(w,h\right)<+\infty\) for \(\nu_{o}\otimes\lambda\)-a.e. \(\left(w,h\right)\in W_{o}\times G\), concluding that the translation \(G\)-space attached to \(\left(W_{o},\nu_{o}\right)\) is totally dissipative.
\end{proof}

\end{document}
